\begin{filecontents*}{paper_draft_reeb_space_sheet_tracking.bib}
@inproceedings{edelsbrunner2008reebspaces,
  author = {Edelsbrunner, Herbert and Harer, John and Patel, Amit K.},
  title = {Reeb Spaces of Piecewise Linear Mappings},
  booktitle = {Proceedings of the Annual Symposium on Computational Geometry},
  pages = {242--250},
  year = {2008},
  doi = {10.1145/1377676.1377720}
}

@article{tierny2017jacobi,
  author = {Tierny, Julien and Carr, Hamish},
  title = {Jacobi Fiber Surfaces for Bivariate Reeb Space Computation},
  journal = {IEEE Transactions on Visualization and Computer Graphics},
  volume = {23},
  number = {1},
  pages = {960--969},
  year = {2017},
  doi = {10.1109/TVCG.2016.2599017}
}

@article{edelsbrunner2008timevarying,
  author = {Edelsbrunner, Herbert and Harer, John and Mascarenhas, Ajith and Pascucci, Valerio and Snoeyink, Jack},
  title = {Time-Varying Reeb Graphs for Continuous Space-Time Data},
  journal = {Computational Geometry},
  volume = {41},
  number = {3},
  pages = {149--166},
  year = {2008},
  doi = {10.1016/j.comgeo.2007.11.001}
}

@article{saikia2017global,
  author = {Saikia, Himangshu and Weinkauf, Tino},
  title = {Global Feature Tracking and Similarity Estimation in Time-Dependent Scalar Fields},
  journal = {Computer Graphics Forum},
  volume = {36},
  number = {3},
  pages = {1--11},
  year = {2017},
  doi = {10.1111/cgf.13163}
}

@article{lukasczyk2020dynamic,
  author = {Lukasczyk, Jonas and Garth, Christoph and Weber, Gunther H. and Biedert, Tim and Maciejewski, Ross and Leitte, Heike},
  title = {Dynamic Nested Tracking Graphs},
  journal = {IEEE Transactions on Visualization and Computer Graphics},
  volume = {26},
  number = {1},
  pages = {249--258},
  year = {2020},
  doi = {10.1109/TVCG.2019.2934368}
}

@misc{sharma2025csp,
  author = {Sharma, Mohit and Masood, Talha Bin and List, Nanna Holmgaard and Hotz, Ingrid and Natarajan, Vijay},
  title = {Continuous Scatterplot and Image Moments for Time-Varying Bivariate Field Analysis of Electronic Structure Evolution},
  year = {2025},
  eprint = {2502.17118},
  archivePrefix = {arXiv},
  primaryClass = {cs.HC},
  doi = {10.48550/arXiv.2502.17118}
}

@misc{hristov2026singular,
  author = {Hristov, Petar and Hotz, Ingrid and Masood, Talha Bin},
  title = {Singular Arrange and Traverse Algorithm for Computing Reeb Spaces of Bivariate PL Maps},
  year = {2026},
  eprint = {2602.21087},
  archivePrefix = {arXiv},
  primaryClass = {cs.CG},
  doi = {10.48550/arXiv.2602.21087}
}
\end{filecontents*}

\documentclass[10pt,journal,compsoc]{IEEEtran}

\usepackage{amsmath}
\usepackage{amssymb}
\usepackage{booktabs}
\usepackage{array}
\usepackage{multirow}
\usepackage{graphicx}
\usepackage{xcolor}
\usepackage{url}
\usepackage[hidelinks]{hyperref}

\newcommand{\R}{\mathbb{R}}
\newcommand{\Domain}{\mathcal{M}}
\newcommand{\Range}{\mathcal{R}}
\newcommand{\Sheets}{\mathcal{S}}
\newcommand{\Sheet}{S}
\newcommand{\Score}{C}
\newcommand{\placeholder}[2]{%
  \begin{center}
  \fbox{%
    \begin{minipage}[c][#1][c]{0.92\linewidth}
      \centering
      #2
    \end{minipage}}
  \end{center}
}
\newcommand{\wideplaceholder}[2]{%
  \begin{center}
  \fbox{%
    \begin{minipage}[c][#1][c]{0.96\textwidth}
      \centering
      #2
    \end{minipage}}
  \end{center}
}

\begin{document}

\title{Tracking Reeb-Space Sheets in Time-Varying Bivariate Fields}

\author{Author~Names~Omitted~for~Draft%
\IEEEcompsocitemizethanks{
\IEEEcompsocthanksitem This is a working draft. Author names, affiliations,
acknowledgments, and final related work will be added later.}}

\markboth{IEEE Transactions on Visualization and Computer Graphics, Draft}%
{Tracking Reeb-Space Sheets in Time-Varying Bivariate Fields}

\IEEEtitleabstractindextext{%
\begin{abstract}
Time-varying bivariate fields arise in many scientific applications, but their
feature-level temporal analysis remains difficult because the features are not
single scalar extrema or contours. Existing visual analysis approaches based on
continuous scatterplots and image moments provide compact summaries of the
overall bivariate distribution, while Reeb-space methods provide a topological
description of a static bivariate map. We connect these two directions by
tracking \emph{Reeb-space sheets} across time. For each timestep, we compute the
Reeb space of two scalar fields, extract the most prominent sheets, and compare
candidate sheet correspondences between adjacent timesteps using complementary
range-space and domain-space measures. Range-space measures compare sheet
footprints using mask intersection-over-union, bounding boxes, area ratios, and
centroid proximity. Domain-space measures compare the mesh vertices associated
with each sheet. The resulting temporal correspondences are visualized using a
Sankey-style view linked with sheet renderings, fiber-surface renderings, and
diagnostic plots. We evaluate the approach on two synthetic torus datasets and
three time-varying molecular electronic-structure datasets represented by hole
and particle natural transition orbital fields. The results show that
range-space sheet overlap gives stable correspondences in the molecular
datasets, while domain overlap exposes a different notion of spatial support
continuity. Together, these measures reveal persistent sheets, abrupt
reconfiguration intervals, and metric disagreements that are difficult to
identify from global bivariate summaries alone.
\end{abstract}

\begin{IEEEkeywords}
Bivariate field, Reeb space, sheet tracking, feature tracking, continuous
scatterplot, fiber surface, time-varying data, molecular visualization.
\end{IEEEkeywords}}

\maketitle

\section{Introduction}

Many scientific simulations produce multiple scalar fields over a common spatial
domain. In electronic-structure analysis, for example, a photoexcited molecule
can be described by hole and particle natural transition orbital (NTO) fields.
The pair of fields identifies where charge is depleted and accumulated during
an electronic transition. Since molecular geometry changes over time, analysts
must compare a sequence of bivariate fields rather than a single pair of
volumes.

One common strategy is to visualize isosurfaces or fiber surfaces at selected
field values. This gives high-quality local evidence, but it does not by itself
provide a compact temporal representation. Another strategy is to summarize each
timestep using a continuous scatterplot (CSP) and then track descriptors of the
CSP over time. This gives an effective overview of global bivariate behavior
and has been used for electronic-structure evolution~\cite{sharma2025csp}. The
limitation is that a global descriptor does not directly identify which
individual bivariate features persist, split, disappear, or reconfigure.

We study a complementary representation based on Reeb spaces. For a bivariate
map \(F=(f,g):\Domain\rightarrow\R^2\), the Reeb space groups domain points
that belong to the same connected component of a fiber \(F^{-1}(a,b)\). In
three-dimensional domains, the two-dimensional cells of the Reeb space can be
interpreted as sheets in the range of the bivariate map. These sheets summarize
regions of coherent bivariate behavior: they are neither individual isosurfaces
nor individual domain regions, but range-space features associated with
connected preimage structure.

This paper proposes a practical pipeline for tracking prominent Reeb-space
sheets in time-varying bivariate data. For each timestep, we compute the Reeb
space, extract the top sheets by area, render each sheet, compute candidate
correspondences between adjacent timesteps, and visualize the resulting temporal
graph using a Sankey-style interface. The method deliberately keeps multiple
similarity notions available. Range-space similarity answers whether two sheets
occupy similar regions in the bivariate range. Domain-space similarity answers
whether two sheets are supported by similar mesh vertices. These two questions
are related but not equivalent, and our experiments show that they often select
different correspondences.

\begin{figure*}[t]
  \wideplaceholder{1.55in}{Figure placeholder: pipeline overview. Show input
  time-varying bivariate fields, Reeb-space computation, top-sheet extraction,
  sheet and fiber-surface rendering, pairwise sheet matching, event analysis,
  and Sankey-style visual exploration.}
  \caption{Overview of the proposed Reeb-space sheet tracking pipeline.}
  \label{fig:pipeline}
\end{figure*}

The main contributions are:
\begin{itemize}
  \item A feature-level temporal representation for time-varying bivariate
  fields based on Reeb-space sheet correspondences.
  \item A set of complementary sheet similarity measures that combine
  range-space shape overlap with domain-space support overlap.
  \item A visual analysis workflow that links temporal Sankey diagrams with
  sheet renderings, fiber-surface renderings, centroid coloring, and
  diagnostic summaries.
  \item A quantitative and qualitative evaluation on synthetic and molecular
  datasets, including metric-agreement analysis and event-interval detection.
\end{itemize}

This draft intentionally leaves the full Related Work section for later. We
still cite directly relevant work where it is necessary to define the method,
position the contribution, or compare capabilities.

\section{Background and Definitions}

\subsection{Bivariate Field}

Let \(\Domain\) be a triangulated spatial domain, typically a tetrahedral or
regular-grid volume. A bivariate field is a continuous or piecewise-linear map
\[
F=(f,g):\Domain\rightarrow\R^2,
\]
where \(f\) and \(g\) are scalar fields sampled on the same domain. In the
molecular datasets used in this paper, \(f\) and \(g\) are NTO-derived scalar
fields, currently configured in the implementation as \texttt{orb00} and
\texttt{orb01}.

The \emph{range} of \(F\) is the subset of \(\R^2\) occupied by field-value
pairs. Each point \(y=(a,b)\in\R^2\) defines a \emph{fiber}
\[
F^{-1}(y)=\{x\in\Domain \mid f(x)=a,\ g(x)=b\}.
\]
For a three-dimensional domain and a bivariate map, a generic fiber is a curve.
A fiber surface is the preimage of a curve in range space, for example the
preimage of \(f=a\), \(f=-a\), \(g=b\), or \(g=-b\). Such surfaces are useful
for interpreting how a selected range-space feature maps back to the spatial
domain.

\subsection{Continuous Scatterplot}

A continuous scatterplot represents the density or measure of domain points
that map to each region of the bivariate range. CSPs are useful for summarizing
the global distribution of field-value pairs. The CSP and image-moment approach
summarizes each timestep by moment descriptors and analyzes the resulting
descriptor sequence in a low-dimensional plot~\cite{sharma2025csp}. Our method
is complementary: it tracks individual Reeb-space sheets rather than only the
global distribution.

\subsection{Reeb Space}

The Reeb space generalizes the Reeb graph from scalar functions to multivariate
maps~\cite{edelsbrunner2008reebspaces}. Define an equivalence relation on
\(\Domain\) by
\[
x\sim x' \quad \text{if} \quad F(x)=F(x') \text{ and } x,x'
\text{ lie in the same connected component of } F^{-1}(F(x)).
\]
The Reeb space is the quotient
\[
\mathcal{R}_F = \Domain/\sim.
\]
For a bivariate map over a three-dimensional domain, the Reeb space is commonly
represented by two-dimensional sheets separated by singular structure. Existing
algorithms compute such structures using Jacobi sets and fiber
surfaces~\cite{tierny2017jacobi}, or arrange-and-traverse style computations
for bivariate piecewise-linear maps~\cite{hristov2026singular}.

\subsection{Sheet}

In this paper, a \emph{sheet} is a two-dimensional cell or connected sheet-like
component in the computed Reeb-space representation. Each sheet has:
\begin{itemize}
  \item a range-space footprint, represented geometrically as polygons in the
  \(f\)-\(g\) plane;
  \item an area in range space;
  \item a centroid and bounding box in range space;
  \item a set of regular domain vertices stored in the sheet information file,
  when such vertices are available.
\end{itemize}
The current pipeline ranks sheets by range-space area and keeps the top
\(N=20\) sheets at each timestep. Zero-regular-vertex sheets are retained
instead of discarded because they can still represent meaningful range-space
geometry. For vertex-overlap metrics, their overlap is treated as zero.

\subsection{Sheet Tracking Problem}

Given a sequence of bivariate fields
\[
F_0,F_1,\ldots,F_T,
\]
and a set of extracted sheets \(\Sheets_i\) for each timestep \(i\), the goal
is to estimate candidate temporal correspondences between sheets in
\(\Sheets_i\) and sheets in \(\Sheets_{i+1}\). We do not assume a one-to-one
matching because sheets can split, merge, appear, disappear, or undergo large
shape changes. Instead, the pipeline stores a weighted bipartite graph for each
adjacent timestep pair.

\section{Method}

\subsection{Input and Preprocessing}

The input is a directory of time-ordered \texttt{.vtu} files. Each file stores
the two scalar arrays used as the bivariate field. The default field names are
configured globally so that all calls to the Reeb-space executable use the same
field pair. The implementation currently uses:
\[
f=\texttt{orb00}, \qquad g=\texttt{orb01}.
\]

For each timestep, the Reeb space is computed using the external Reeb-space
program. The pipeline stores two primary outputs:
\begin{itemize}
  \item \texttt{.rs}: the Reeb-space geometry and sheet structure used for
  rendering and geometry export;
  \item \texttt{.rsi}: sheet information used by the Python pipeline, including
  sheet identifiers, area/rank data, and associated regular vertices.
\end{itemize}
Only timesteps with usable \texttt{.rs} and \texttt{.rsi} files are included in
the main tracking graph. In the main analysis branch, the pipeline does not use
perturbed inputs for matching; degenerate or failed timesteps are excluded from
the downstream analysis unless a valid sheet representation is present.

\subsection{Top-Sheet Selection}

Let \(\Sheets_i^\mathrm{all}\) be the set of sheets at timestep \(i\). We sort
the sheets by range-space area and keep the top \(N\) sheets:
\[
\Sheets_i = \operatorname{TopN}(\Sheets_i^\mathrm{all}, N).
\]
The experiments use \(N=20\). This keeps the temporal graph readable and
focuses the analysis on dominant bivariate features. The choice of \(N\) is a
parameter and can be varied in future sensitivity studies.

\subsection{Range-Space Geometry Export}

For each selected sheet, the pipeline exports a polygonal representation of the
sheet footprint in range space. The exported VTP cache stores a collection of
sheet polygons for a timestep, not a separate VTP file per sheet. These cached
VTP files are used for shape comparison and sheet rendering. To ensure visual
comparability across time, sheet images are rendered using global bounds and a
fixed image size. This avoids a misleading effect where sheets appear to have
similar size because each image is independently zoomed.

\begin{figure}[t]
  \placeholder{1.35in}{Figure placeholder: one timestep with the full
  Reeb-space sheet collection, top sheets highlighted, and a zoomed example of
  one sheet footprint in range space.}
  \caption{A Reeb-space sheet is represented by its footprint in the bivariate
  range and by its associated domain support.}
  \label{fig:sheet-definition}
\end{figure}

\subsection{Sheet Rendering}

Each selected sheet is rendered as an image. The rendering stage uses a uniform
sheet color to avoid encoding unrelated orbital or metric information into the
sheet image itself. A faint context can be rendered from the full sheet set,
while the selected sheet remains visually dominant. The same global range-space
camera and image size are used for all timesteps, so changes in sheet size and
position can be compared visually.

\subsection{Fiber-Surface Rendering}

For selected sheet and timestep combinations, the pipeline renders
fiber-surface images using a ParaView state file. The isovalues are configured
globally. For a value \(a\) of field \(f\) and a value \(b\) of field \(g\),
the pipeline renders surfaces corresponding to:
\[
f=+a,\quad f=-a,\quad g=+b,\quad g=-b.
\]
The Reeb-space executable currently generates fiber surfaces for the complete
data. The Python stage then thresholds or filters those surfaces to the top
sheets of interest and renders the corresponding image. If a surface does not
intersect the selected sheet, the generated VTP may have zero points and zero
cells; such cases are retained as empty evidence rather than treated as a
pipeline error.

\subsection{Range-Space Sheet Similarity}

For each adjacent timestep pair \((i,i+1)\), every source sheet
\(\Sheet\in\Sheets_i\) is compared with every target sheet
\(\Sheet'\in\Sheets_{i+1}\). The first group of metrics compares sheet
geometry in the bivariate range.

\subsubsection{Shape IoU}

The range-space polygons of each sheet are rasterized into masks over global
range-space bounds. The shape intersection-over-union is
\[
\operatorname{shapeIoU}(\Sheet,\Sheet') =
\frac{|M_\Sheet\cap M_{\Sheet'}|}{|M_\Sheet\cup M_{\Sheet'}|},
\]
where \(M_\Sheet\) and \(M_{\Sheet'}\) are binary raster masks.

\subsubsection{Bounding-Box IoU}

Let \(B_\Sheet\) and \(B_{\Sheet'}\) be axis-aligned bounding boxes in range
space. The bounding-box IoU is
\[
\operatorname{bboxIoU}(\Sheet,\Sheet') =
\frac{\operatorname{area}(B_\Sheet\cap B_{\Sheet'})}
{\operatorname{area}(B_\Sheet\cup B_{\Sheet'})}.
\]
This is cheaper and coarser than mask-based shape IoU.

\subsubsection{Area Ratio}

The area-ratio score measures whether two sheets have comparable area:
\[
\operatorname{areaRatio}(\Sheet,\Sheet') =
\frac{\min(A_\Sheet,A_{\Sheet'})}{\max(A_\Sheet,A_{\Sheet'})}.
\]
It does not encode spatial location in the range plane.

\subsubsection{Centroid Similarity}

Let \(c_\Sheet\) and \(c_{\Sheet'}\) be sheet centroids, and let \(D\) be the
diagonal of the global range-space bounds. The centroid similarity is
\[
\operatorname{centroidSim}(\Sheet,\Sheet') =
\exp\left(-\frac{\|c_\Sheet-c_{\Sheet'}\|}{D}\right).
\]
This metric is useful as a weak positional cue but is not sufficiently
discriminative as a standalone matching score.

\subsection{Domain-Space Sheet Similarity}

The RSI file stores regular domain vertices associated with each sheet. Let
\(V_\Sheet\) and \(V_{\Sheet'}\) be the corresponding vertex sets. We compute
the support Jaccard score:
\[
\operatorname{supportJaccard}(\Sheet,\Sheet') =
\frac{|V_\Sheet\cap V_{\Sheet'}|}{|V_\Sheet\cup V_{\Sheet'}|}.
\]
If either sheet has no stored regular vertices, the score is zero. This choice
is conservative: it avoids inventing vertex overlap where none is recorded.
Importantly, a zero support score does not imply that the sheet has no
range-space geometry.

\subsection{Combined Matching Score}

The default combined sheet score is a weighted sum:
\[
\begin{aligned}
\Score(\Sheet,\Sheet') ={}&
0.40\,\operatorname{shapeIoU}(\Sheet,\Sheet') \\
&+0.30\,\operatorname{supportJaccard}(\Sheet,\Sheet') \\
&+0.15\,\operatorname{areaRatio}(\Sheet,\Sheet') \\
&+0.10\,\operatorname{bboxIoU}(\Sheet,\Sheet') \\
&+0.05\,\operatorname{centroidSim}(\Sheet,\Sheet').
\end{aligned}
\]
The weights are not claimed to be universal. They are current defaults chosen
to privilege range-shape agreement while retaining domain support and coarse
geometric cues. The analysis stage reports individual metric behavior so that
users can understand when the combined score is driven by range shape, domain
support, or size/location effects.

\subsection{Temporal Graph Construction}

For every adjacent timestep pair, the method constructs a complete candidate
bipartite graph between selected sheets:
\[
G_i=(\Sheets_i,\Sheets_{i+1},E_i),
\]
where each edge \(e=(\Sheet,\Sheet')\) stores all metric values and the
combined score. The final temporal graph is a sequence of these adjacent
bipartite graphs. The viewer can display the graph using the combined score,
domain overlap, range overlap, or a hybrid score.

The current method does not solve a global one-to-one assignment. This is
intentional. Sheet evolution can be many-to-many, and retaining all candidate
edges supports visual inspection of split and merge hypotheses.

\subsection{Event Diagnostics}

To identify intervals that deserve closer inspection, we compute an event score
for each adjacent timestep pair and for several thresholds
\(\theta\in\{0.3,0.4,0.5,0.6,0.7\}\). For a source sheet, define
\[
\operatorname{bestOut}(\Sheet)=
\max_{\Sheet'\in\Sheets_{i+1}}\Score(\Sheet,\Sheet').
\]
For a target sheet, define
\[
\operatorname{bestIn}(\Sheet')=
\max_{\Sheet\in\Sheets_i}\Score(\Sheet,\Sheet').
\]
At threshold \(\theta\), we count:
\begin{itemize}
  \item weak outgoing sheets: \(\operatorname{bestOut}(\Sheet)<\theta\);
  \item weak incoming sheets: \(\operatorname{bestIn}(\Sheet')<\theta\);
  \item possible splits: one source has two or more targets with score
  at least \(\theta\);
  \item possible merges: one target has two or more sources with score
  at least \(\theta\).
\end{itemize}
The exploratory event score is
\[
\begin{aligned}
E_i(\theta)={}&
W^\mathrm{out}_i + W^\mathrm{in}_i
+ \frac{1}{2}(P^\mathrm{split}_i + P^\mathrm{merge}_i) \\
&+ (1-\overline{B^\mathrm{out}_i})|\Sheets_i|,
\end{aligned}
\]
where \(\overline{B^\mathrm{out}_i}\) is the mean best outgoing combined score.
Large values indicate intervals with weak continuations, many potential
reconfigurations, or generally low sheet similarity. The threshold
\(\theta=0.5\) is used as the default for ranked intervals, while all threshold
values are exported for sensitivity analysis.

\subsection{Sheet Lifetimes}

For a given threshold, we compute a greedy lifetime summary. Starting from a
sheet at timestep \(i\), the analysis follows the best outgoing edge while its
score remains above threshold. The resulting chain length is a diagnostic for
persistence. It is not a proof of a unique physical track, but it helps
identify long-lived prominent sheets and intervals where tracks fragment.

\subsection{Visual Encoding and Interaction}

The unified viewer displays sheets as nodes and candidate correspondences as
links. Node color can encode rank, metric value, or centroid position. For
centroid-based coloring, global symmetric range-space extents are used so that
colors are comparable across timesteps. Two centroid color maps are currently
available: a four-corner interpolation map and an axis-oriented map in which
distance from the origin increases saturation along red-blue directions.

Clicking or hovering a node shows the corresponding sheet image and
fiber-surface images. Clicking or hovering a link opens a side-by-side
comparison of the two sheets and their fiber-surface images. The image detail
view supports linked pan and zoom so that both sides of a comparison can be
inspected at the same scale.

\begin{figure*}[t]
  \wideplaceholder{1.75in}{Figure placeholder: unified viewer screenshot.
  Include a Sankey diagram with node/link coloring, a selected link, side-by-side
  sheet images, and side-by-side fiber-surface images with linked zoom.}
  \caption{The visual interface links temporal sheet correspondences with
  range-space sheet images and spatial fiber-surface evidence.}
  \label{fig:viewer}
\end{figure*}

\section{Implementation}

The pipeline is implemented as a sequence of Python stages. The stages are:
\begin{enumerate}
  \item compute Reeb-space files from VTU inputs;
  \item convert RSI files into JSON used by downstream stages;
  \item compute range-space sheet shape matches;
  \item compute domain-space vertex overlaps;
  \item render sheet images;
  \item render sheet-restricted fiber-surface images;
  \item build the unified Sankey viewer;
  \item run tracking-analysis diagnostics.
\end{enumerate}
The current stage naming in the repository uses the sequence
1, 2, 3A, 3B, 4A, 4B, 5, with the paper-analysis stage added afterward as an
optional diagnostic stage.

The Reeb-space computation is delegated to the external \texttt{fv99}
executable. The Python pipeline passes the configured field names and uses the
resulting \texttt{.rs} and \texttt{.rsi} files for all downstream analysis. VTP
exports are cached because sheet geometry extraction can be expensive. Cache
reuse is conservative: cached geometry is reused unless the user explicitly
requests a rebuild or clean cache.

The tracking-analysis stage consumes the existing unified-viewer JSON and
writes CSV, JSON, PNG, and PDF outputs. It does not recompute Reeb spaces,
sheet VTPs, sheet images, or fiber-surface images. This separation keeps the
paper diagnostics reproducible from a fixed set of generated viewer data.

\section{Datasets}

We evaluate on five datasets: three molecular electronic-structure sequences
and two synthetic torus sequences. The molecular datasets are named
\texttt{MVK\_s1}, \texttt{MVK\_s2}, and \texttt{stilbene}. The synthetic
datasets are \texttt{torusGrowing} and \texttt{torusMoving}. The torus datasets
are used for controlled behavior: one changes feature scale over time, and the
other moves the feature center smoothly over time.

\begin{table}[t]
\caption{Datasets used in the current analysis. Each timestep keeps the top 20
sheets by range-space area. Zero-vertex sheets are retained.}
\label{tab:datasets}
\centering
\begin{tabular}{lrrr}
\toprule
Dataset & Timesteps & Top sheets & Zero-vertex sheets \\
\midrule
MVK\_s1 & 82 & 1640 & 116 \\
MVK\_s2 & 82 & 1640 & 68 \\
stilbene & 698 & 13960 & 148 \\
torusGrowing & 25 & 500 & 188 \\
torusMoving & 25 & 500 & 151 \\
\bottomrule
\end{tabular}
\end{table}

\begin{figure*}[t]
  \wideplaceholder{1.8in}{Figure placeholder: dataset overview. Show one or
  two representative timesteps from each dataset, including the molecule or
  synthetic volume, the CSP/range-space view, and selected Reeb-space sheets.}
  \caption{Representative timesteps from the molecular and synthetic datasets.}
  \label{fig:datasets}
\end{figure*}

\section{Results}

\subsection{Metric Agreement}

We first ask how often an individual metric selects the same best target as the
default combined score. For each source sheet, we choose the target with maximum
value according to one metric and compare it with the target chosen by the
combined score. Table~\ref{tab:metric-agreement} summarizes the agreement.

\begin{table*}[t]
\caption{Best-target agreement with the default combined score. Values are
percentages of source sheets for which an individual metric selects the same
target as the combined score.}
\label{tab:metric-agreement}
\centering
\begin{tabular}{lrrrrrr}
\toprule
Dataset & Shape IoU & BBox IoU & Support Jaccard & Area ratio & Centroid & Domain overlap \\
\midrule
MVK\_s1 & 95.6 & 94.2 & 50.0 & 65.9 & 90.9 & 9.7 \\
MVK\_s2 & 93.7 & 91.4 & 48.6 & 58.1 & 88.0 & 9.6 \\
stilbene & 96.1 & 89.3 & 40.5 & 47.7 & 87.1 & 8.7 \\
torusGrowing & 85.4 & 44.2 & 70.2 & 43.1 & 42.7 & 16.0 \\
torusMoving & 37.7 & 29.2 & 46.5 & 89.2 & 31.2 & 1.8 \\
\bottomrule
\end{tabular}
\end{table*}

For the molecular datasets, shape IoU is highly consistent with the combined
score. This is partly expected because shape IoU has the largest weight in the
combined score. However, the result is still useful: it shows that the combined
score is largely governed by range-space sheet footprint overlap in these
datasets. Bounding-box IoU is also strong but less precise. Centroid similarity
often agrees with the combined score in the molecular data, but visual
inspection shows that it is a weaker standalone metric because many candidates
can be spatially close in the range plane.

Domain overlap behaves differently. The domain-overlap best target agrees with
the combined shape target in only \(8.7\%\)--\(9.7\%\) of source sheets in the
molecular datasets. This low agreement is not necessarily a failure. It means
that range-shape continuity and domain-support continuity answer different
questions. A sheet can keep a similar range-space footprint while being
supported by different domain vertices, or it can preserve some domain support
while changing its range-space shape.

The synthetic datasets expose different behavior. In \texttt{torusGrowing},
support Jaccard is stronger than in the molecular data, while bounding-box and
centroid metrics are less reliable. In \texttt{torusMoving}, area ratio agrees
with the combined score more often than shape IoU because the dominant feature
changes position more than size. This result supports the need to report
multiple metrics rather than relying on a single universal score.

\begin{figure*}[t]
  \wideplaceholder{1.55in}{Figure placeholder: metric comparison plots. Include
  distributions of shape IoU, support Jaccard, centroid similarity, and combined
  score; include domain-vs-shape best-target disagreement.}
  \caption{Metric diagnostics reveal that range-space and domain-space
  similarities capture complementary notions of sheet continuity.}
  \label{fig:metric-diagnostics}
\end{figure*}

\subsection{Event Intervals}

Table~\ref{tab:event-intervals} reports the top event intervals according to
the default threshold \(\theta=0.5\). These intervals are not final scientific
interpretations; they are candidates for visual inspection and domain expert
analysis.

\begin{table*}[t]
\caption{Highest-ranked event intervals at threshold \(\theta=0.5\). Mean best
score is the mean best outgoing combined score. Weak source and weak target
counts are out of 20 top sheets per timestep.}
\label{tab:event-intervals}
\centering
\begin{tabular}{llrrrr}
\toprule
Dataset & Interval & Event score & Mean best & Weak source & Weak target \\
\midrule
MVK\_s1 & step\_01200 $\rightarrow$ step\_01220 & 48.08 & 0.396 & 18 & 18 \\
MVK\_s1 & step\_01180 $\rightarrow$ step\_01200 & 37.70 & 0.415 & 13 & 13 \\
MVK\_s2 & step\_01140 $\rightarrow$ step\_01160 & 42.14 & 0.393 & 15 & 15 \\
MVK\_s2 & step\_01160 $\rightarrow$ step\_01180 & 40.23 & 0.388 & 14 & 14 \\
stilbene & step\_11915 $\rightarrow$ step\_11920 & 53.57 & 0.321 & 20 & 20 \\
stilbene & step\_11920 $\rightarrow$ step\_11925 & 53.30 & 0.335 & 20 & 20 \\
torusGrowing & step\_00200 $\rightarrow$ step\_00220 & 47.28 & 0.336 & 16 & 16 \\
torusMoving & step\_00200 $\rightarrow$ step\_00220 & 48.21 & 0.289 & 16 & 16 \\
\bottomrule
\end{tabular}
\end{table*}

The molecular datasets show compact event regions. In \texttt{MVK\_s1}, the
strongest region is around \texttt{step\_01180}--\texttt{step\_01220}, with the
largest event score at \texttt{step\_01200} to \texttt{step\_01220}. In
\texttt{MVK\_s2}, the strongest region is around
\texttt{step\_01120}--\texttt{step\_01180}. In \texttt{stilbene}, the strongest
intervals occur around \texttt{step\_11915}--\texttt{step\_11925}, where all
top sheets have weak outgoing and incoming continuations under the default
threshold.

The synthetic datasets show frequent weak tracks, partly because many top
sheets are small or have zero stored regular vertices. These cases are useful
stress tests for the analysis but should not be interpreted in the same way as
the molecular events without further synthetic ground-truth labeling.

\begin{figure*}[t]
  \wideplaceholder{1.6in}{Figure placeholder: event-score curves for MVK\_s1,
  MVK\_s2, and stilbene. Mark the intervals reported in
  Table~\ref{tab:event-intervals}.}
  \caption{Event-score curves help prioritize temporal intervals for detailed
  inspection.}
  \label{fig:event-curves}
\end{figure*}

\subsection{Sheet Lifetimes}

The lifetime diagnostic summarizes how long greedy best-score tracks persist at
threshold \(\theta=0.5\). The maximum lifetime reaches all 82 timesteps for
\texttt{MVK\_s1}, 57 timesteps for \texttt{MVK\_s2}, 139 timesteps for
\texttt{stilbene}, and all 25 timesteps for both synthetic datasets. However,
the median lifetime is short: one timestep for \texttt{MVK\_s1},
\texttt{MVK\_s2}, \texttt{torusGrowing}, and \texttt{torusMoving}, and two
timesteps for \texttt{stilbene}. This indicates a mixture of long-lived
dominant sheets and many short-lived or unstable lower-ranked sheets.

\begin{figure}[t]
  \placeholder{1.25in}{Figure placeholder: lifetime histogram for each dataset
  at threshold 0.5, with insets for the longest tracks.}
  \caption{Greedy sheet lifetimes reveal long-lived dominant sheets amid many
  short-lived top-sheet candidates.}
  \label{fig:lifetime}
\end{figure}

\subsection{Visual Case Studies}

\subsubsection{MVK\_s1}

The strongest event interval occurs around \texttt{step\_01200} to
\texttt{step\_01220}. The mean best combined score is 0.396, and 18 of the 20
top sheets have weak outgoing and incoming continuations at threshold 0.5. This
interval should be shown with side-by-side sheet images and fiber-surface
images. The expected visual evidence is a reconfiguration of dominant
range-space sheets rather than a simple smooth continuation.

\begin{figure*}[t]
  \wideplaceholder{1.8in}{Figure placeholder: MVK\_s1 case study. Show
  \texttt{step\_01180}, \texttt{step\_01200}, and \texttt{step\_01220}; include
  Sankey subgraph, selected sheet images, fiber surfaces, and event-score
  marker.}
  \caption{Candidate event in \texttt{MVK\_s1}.}
  \label{fig:mvk-s1-case}
\end{figure*}

\subsubsection{MVK\_s2}

The main event region occurs around \texttt{step\_01120} to
\texttt{step\_01180}. Two adjacent intervals,
\texttt{step\_01140} to \texttt{step\_01160} and \texttt{step\_01160} to
\texttt{step\_01180}, have high event scores and low mean best continuation.
This suggests a sustained transition rather than a single isolated timestep.

\begin{figure*}[t]
  \wideplaceholder{1.8in}{Figure placeholder: MVK\_s2 case study with temporal
  window \texttt{step\_01120}--\texttt{step\_01180}.}
  \caption{Candidate event region in \texttt{MVK\_s2}.}
  \label{fig:mvk-s2-case}
\end{figure*}

\subsubsection{Stilbene}

The \texttt{stilbene} dataset contains 698 usable timesteps and therefore
benefits most from a temporal overview. The strongest event region appears
around \texttt{step\_11915} to \texttt{step\_11925}. Both adjacent intervals
have all 20 source and target sheets below threshold 0.5. This region should be
examined in the viewer together with molecular geometry and fiber-surface
images.

\begin{figure*}[t]
  \wideplaceholder{1.8in}{Figure placeholder: stilbene case study around
  \texttt{step\_11915}--\texttt{step\_11925}. Include sheet renderings,
  fiber-surface renderings, and linked temporal context.}
  \caption{Candidate event region in \texttt{stilbene}.}
  \label{fig:stilbene-case}
\end{figure*}

\subsubsection{Synthetic Torus}

The synthetic torus datasets are used to test controlled changes. The
\texttt{torusGrowing} sequence should emphasize scale change, while
\texttt{torusMoving} should emphasize translation in the range-space
representation. Current results show that metric behavior changes substantially
between these two cases. In the moving sequence, area ratio remains high while
shape IoU and centroid metrics may disagree with the combined score, depending
on the exact geometry and sheet ranking. This makes the synthetic datasets
useful for validating metric sensitivity.

\begin{figure*}[t]
  \wideplaceholder{1.6in}{Figure placeholder: synthetic torus results. Compare
  torusGrowing and torusMoving with known expected behavior, selected sheet
  tracks, and metric curves.}
  \caption{Synthetic datasets provide controlled tests for scale and motion.}
  \label{fig:torus-case}
\end{figure*}

\section{Comparison With Existing Approaches}

The goal of this section is not to replace a full Related Work section. Instead
it summarizes how the proposed pipeline differs from several directly relevant
families of methods.

\begin{table*}[t]
\caption{Capability comparison with representative prior approaches. The table
is qualitative and should be refined after the full Related Work section is
written.}
\label{tab:comparison}
\centering
\begin{tabular}{p{0.23\textwidth}p{0.14\textwidth}p{0.14\textwidth}p{0.14\textwidth}p{0.14\textwidth}p{0.14\textwidth}}
\toprule
Approach & Bivariate input & Reeb-space features & Temporal tracking & Domain/range metrics & Linked feature images \\
\midrule
CSP image moments~\cite{sharma2025csp} & Yes & No & Descriptor curve & Global descriptors & Fiber surfaces for selected states \\
Static Reeb-space computation~\cite{tierny2017jacobi,hristov2026singular} & Yes & Yes & No & Static topology & Not primary focus \\
Time-varying Reeb graphs~\cite{edelsbrunner2008timevarying} & Scalar & Reeb graph & Yes & Topological graph evolution & Level-set extraction support \\
Scalar feature tracking~\cite{saikia2017global,lukasczyk2020dynamic} & Scalar or component based & Merge-tree/component features & Yes & Overlap and similarity & Application dependent \\
Ours & Yes & Reeb-space sheets & Yes & Domain and range metrics & Sheet and fiber-surface views \\
\bottomrule
\end{tabular}
\end{table*}

Compared with CSP image moments, our method is less compact but more
feature-specific. A CSP moment curve can quickly identify unusual timesteps, but
it does not directly indicate which Reeb-space sheet persisted or changed.
Compared with static Reeb-space computation, our method adds temporal
correspondence and visual diagnostics. Compared with scalar feature tracking,
our method operates on bivariate Reeb-space sheets rather than scalar
components or merge-tree regions.

\section{Discussion}

\subsection{Why Domain and Range Metrics Both Matter}

The low agreement between domain-overlap and range-shape best targets is one of
the most important observations. It shows that ``same feature'' is not a single
unambiguous concept for bivariate Reeb-space sheets. If the scientific question
concerns a stable relationship between the two scalar fields, range-space
similarity is appropriate. If the question concerns persistence of a spatial
region in the original domain, domain support overlap is more relevant. The
viewer therefore exposes both notions instead of hiding them inside one score.

\subsection{Interpreting Zero-Vertex Sheets}

Some high-rank sheets have zero stored regular vertices. The pipeline now keeps
these sheets because their range-space footprint and sheet area can still be
meaningful. Vertex-based overlap for such sheets is zero by construction. In
the paper, these sheets should be discussed carefully: they are valid
range-space entities, but they provide no regular-vertex support evidence under
the current RSI representation.

\subsection{Threshold Sensitivity}

The event threshold \(\theta=0.5\) is a useful exploratory setting, but it is
not a calibrated physical threshold. The analysis stage exports event scores
and lifetimes for thresholds from 0.3 to 0.7. A final paper should include a
threshold-sensitivity figure showing which event intervals remain stable across
thresholds. Intervals that appear only at one threshold should be treated as
weak evidence.

\subsection{Computational Considerations}

The most expensive parts of the pipeline are Reeb-space computation, sheet
geometry export, sheet rendering, and fiber-surface rendering. Caching reduces
recomputation, but cache consistency is important: all stages should use the
same input VTU, field names, and Reeb-space outputs. Rendering uses global
sheet bounds for comparability, which can increase memory use because a larger
context must be available for each image.

\section{Limitations}

The current method has several limitations.

First, the matching score is heuristic. The weights are practical defaults, not
learned or theoretically optimal. We report individual metric behavior to make
this visible, but a final system should allow users to tune weights and inspect
sensitivity.

Second, the current temporal graph stores adjacent-pair candidate matches. It
does not solve a global optimization problem over the full time sequence. This
is appropriate for exploratory visual analysis, but it can produce fragmented
tracks when short-term scores are ambiguous.

Third, the method depends on successful Reeb-space computation and geometry
export. Degenerate bivariate inputs can cause the external computation to fail.
The main branch of the analysis avoids perturbation-based correction so that
all reported results correspond to unperturbed input data.

Fourth, the current RSI representation stores regular vertices but not all
singular-vertex membership information per sheet. This limits the precision of
domain-support comparisons for sheets whose evidence is concentrated near
singular structure.

Fifth, the event score is an analysis heuristic. It is useful for prioritizing
inspection but should not be presented as a mathematically complete event
detector without additional validation.

Finally, the synthetic datasets currently provide controlled qualitative tests
but not a full ground-truth benchmark. A stronger evaluation would include
known sheet correspondences, controlled split/merge events, and quantitative
accuracy against ground truth.

\section{Future Work}

Future work should address four directions. First, we plan to add a formal
threshold and weight sensitivity study. Second, we will improve the RSI
representation by storing singular-vertex membership per sheet, enabling richer
domain-overlap metrics. Third, we will design synthetic benchmarks with known
sheet correspondences and controlled degeneracies. Fourth, we will add global
tracking algorithms that use the adjacent-pair graph as input but optimize
tracks over the full sequence.

\section{Conclusion}

We presented a Reeb-space sheet tracking pipeline for time-varying bivariate
fields. The method extracts prominent Reeb-space sheets, compares them across
adjacent timesteps using range-space and domain-space metrics, and visualizes
the resulting correspondences with linked sheet and fiber-surface evidence.
Experiments on molecular and synthetic datasets show that range-space shape
overlap provides stable correspondences in molecular sequences, while
domain-overlap metrics expose complementary spatial-support behavior. The
approach bridges global bivariate summaries and feature-level temporal
inspection, supporting the identification of persistent sheets and important
transition intervals in time-varying bivariate data.

\section*{Acknowledgments}

Acknowledgments will be added in the final version.

\bibliographystyle{IEEEtran}
\bibliography{paper_draft_reeb_space_sheet_tracking}

\end{document}
